\subsection{同态与同伦的复形}

命题3. --- 设 C， \(\tilde{C}\) , \(C'\) , \(\tilde{C}'\) 为四个 A-复形， \(u: \tilde{C} \to C\) , \(v: \tilde{C} \to C\) , \(u': C' \to \tilde{C}'\) 和 \(v': C' \to \tilde{C}'\) 为四个复形态射。

\begin{enumerate}
\def\labelenumi{\alph{enumi})}
\item
  如果 u 和 \(u'\) 分别与 v 和 \(v'\) 同伦，则这两个态射 \(\operatorname{Homgr}_{\mathbf{A}}(u, u')\) 和 \(\operatorname{Homgr}_{\mathbf{A}}(v, v')\) （从 \(\operatorname{Homgr}_{\mathbf{A}}(\mathbf{C}, \mathbf{C}')\) 到 \(\operatorname{Homgr}_{\mathbf{A}}(\tilde{\mathbf{C}}, \tilde{\mathbf{C}}')\) ）是同伦的。
\item
  如果 u 和 u' 是同伦等价，Homgr\_A(u, u') 是一个同伦等价。
\item
  如果 C 或 C' 同伦于零，则 Homgr\_A(C, C') 同伦于零。
\end{enumerate}

我们用同一个字母 \(d\) 表示复形 C, C \(_{1}\) ，C \(^{\prime}\) ，C \(_{1}^{\prime}\) 的微分，并用 D 表示 Homgr \(_{A}\) (C, C') 和 Homgr \(_{A}\) (C \(_{1}\) ，C \(_{1}^{\prime}\) ) 的微分。若 \(u\) （相应地， \(u^{\prime}\) ）同伦于 \(v\) （相应地， \(v^{\prime}\) ），则存在一个次数为 1 的分次同态，

\[
w: \mathrm{C} _ {1} \rightarrow \mathrm{C} \quad (\text { resp. } w ^ {\prime}: \mathrm{C} ^ {\prime} \rightarrow \mathrm{C} _ {1} ^ {\prime})
\]

使得

\[
u - v = d w + w d \quad (\text { resp. } u ^ {\prime} - v ^ {\prime} = d w ^ {\prime} + w ^ {\prime}\tag{2}
\]

设 W : Homgr \(_{A}\) (C, C') → Homgr \(_{A}\) (C \(_{1}\) ，C \(_{1}\) ) 是次数为 1 的分次同态，使得对于 \(f \in\) Homgr \(_{A}\) (C, C') \(_{n}\) , \(n \in\) Z，我们有

\[
\mathrm{W} (f) = w ^ {\prime} f u + (- 1) ^ {n} v ^ {\prime} f w.\tag{3}
\]

于是我们有

\[
\begin{array}{r l} (\mathrm{DW} + \mathrm{WD}) (f) & = \mathrm{D} [ w ^ {\prime} f u + (- 1) ^ {n} v ^ {\prime} f w ] + \mathrm{W} [ d f - (- 1) ^ {n} f d ] \\ & = d w ^ {\prime} f u - (- 1) ^ {n + 1} w ^ {\prime} f u d + (- 1) ^ {n} d v ^ {\prime} f w + v ^ {\prime} f w d \\ & \quad + w ^ {\prime} d f u + (- 1) ^ {n + 1} v ^ {\prime} d f w - (- 1) ^ {n} w ^ {\prime} f d u + v ^ {\prime} f d w \\ & = (d w ^ {\prime} + w ^ {\prime} d) f u + v ^ {\prime} f (w d + d w) \\ & = (u ^ {\prime} - v ^ {\prime}) f u + v ^ {\prime} f (u - v) = u ^ {\prime} f u - v ^ {\prime} f v; \end{array}
\]

即 DW + WD = Homgr\_A (u, u')-Homgr\_A (v, v')，由此得 a)。

让我们证明 b)。如果 u 和 \(u'\) 是同伦等价，令 \(\alpha: C \to \tilde{C}\) 和 \(\alpha': \tilde{C}' \to C'\) 是复形的态射，使得 \(u \circ \alpha\) , \(\alpha \circ u\) , \(u' \circ \alpha'\) , \(\alpha' \circ u'\) 分别同伦于 \(Id_{C}\) , \(Id_{\tilde{C}}\) , \(Id_{\tilde{C}'}\) , \(Id_{C'}\) 。则 Homgr \((u, u') \circ \text{Homgr } (\alpha, \alpha')\) ，它等于 Homgr \((\alpha \circ u, u' \circ \alpha')\) ，通过a）同伦于

\[
\operatorname{Homgr} _ {A} \left(\operatorname{Id} _ {C}, \operatorname{Id} _ {C ^ {\prime}}\right) = \operatorname{Id} _ {\operatorname{Homgr} (C, C ^ {\prime})}:
\]

同样地 Homgr \((\alpha, \alpha') \circ\) Homgr \((u, u')\) 同伦于 \(\operatorname{Id}_{\operatorname{Homgr}(\tilde{\mathbb{C}}, \tilde{\mathbb{C}})}\) ，从而 \(b\) ).

最后 c) 由 b) 推出（应用于 \(\tilde{C}\) 或 \(\tilde{C}'\) 为零的情形）。

推论 1. --- 若 C 是分裂的且 H(C) 是投射的（分别地，若 C' 是分裂的且对每个 n，H \(_{n}\) (C') 是内射的），则典范同态

\[
\lambda (\mathrm{C}, \mathrm{C} ^ {\prime}): \mathrm{H} (\operatorname{Homgr} _ {\mathrm{A}} (\mathrm{C}, \mathrm{C} ^ {\prime})) \rightarrow \operatorname{Homgr} _ {\mathrm{A}} (\mathrm{H} (\mathrm{C}), \mathrm{H} (\mathrm{C} ^ {\prime}))
\]

是双射。

例如假设 C' 是分裂的且对每个 n，H(C') 是内射的，而 C 是分裂的且 H(C) 是投射的情形可类似地证明。由 X，第 35 页，定义 6，存在一个同伦等价 \(u': C' \to H(C')\) ；由命题 3，Homgr\_A (1, \(u'\) ) 是从 Homgr\_A (C, C') 到 Homgr\_A (C, H(C')) 的同伦等价；由于

\[
\operatorname{Homgr} _ {\mathbf {A}} \left(1, \mathrm{H} \left(u ^ {\prime}\right)\right) \circ \lambda (\mathrm{C}, \mathrm{C} ^ {\prime}) = \lambda (\mathrm{C}, \mathrm{H} \left(\mathrm{C} ^ {\prime}\right)) \circ \mathrm{H} \left(\operatorname{Homgr} _ {\mathbf {A}} \left(1, u ^ {\prime}\right)\right)
\]

并且 \(\operatorname{Homgr}_{\mathbf{A}}(1, \mathrm{H}(u'))\) 和 \(\operatorname{H}(\operatorname{Homgr}_{\mathbf{A}}(1, u'))\) 是双射的，我们只需证明 \(\lambda(\mathbf{C},\mathbf{H}(\mathbf{C}^{\prime}))\) 是双射的，这使我们回到以下情形： \(\mathbf{C}'\) 是内射的且具有零微分。

那么典范正合序列（X，第28页）

\[
0 \rightarrow \mathrm{B} (\mathrm{C}) \xrightarrow {i} \mathrm{C} \xrightarrow {p} \mathrm{C} / \mathrm{B} (\mathrm{C}) \rightarrow 0\tag{III}
\]

\[
0 \rightarrow \mathrm{H} (\mathrm{C}) \stackrel {{j}} {{\rightarrow}} \mathrm{C} / \mathrm{B} (\mathrm{C}) \stackrel {{\delta}} {{\rightarrow}} \mathrm{B} (\mathrm{C}) \rightarrow 0\tag{IV}
\]

给出正合序列（X，第83页，命题2，a））

\[
0 \rightarrow \operatorname{Homgr} _ {\mathrm{A}} (\mathrm{C} / \mathrm{B} (\mathrm{C}), \mathrm{C} ^ {\prime}) \xrightarrow {\mathrm{P}} \operatorname{Homgr} _ {\mathrm{A}} (\mathrm{C}, \mathrm{C} ^ {\prime}) \xrightarrow {\mathrm{I}} \operatorname{Homgr} _ {\mathrm{A}} (\mathrm{B} (\mathrm{C}), \mathrm{C} ^ {\prime}) \rightarrow 0
\]

\[
0 \rightarrow \operatorname{Homgr} _ {\mathrm{A}} (\mathrm{B} (\mathrm{C}), \mathrm{C} ^ {\prime}) \stackrel {{\Delta}} {{\rightarrow}} \operatorname{Homgr} _ {\mathrm{A}} (\mathrm{C} / \mathrm{B} (\mathrm{C}), \mathrm{C} ^ {\prime}) \stackrel {{\mathrm{J}}} {{\rightarrow}} \operatorname{Homgr} _ {\mathrm{A}} (\mathrm{H} (\mathrm{C}), \mathrm{C} ^ {\prime}) \rightarrow 0.
\]

由于 \(d_{\mathbf{C}} = i \circ \delta \circ p\) ，Homgr(C, C') 的微分 D 由下式给出 \(\mathrm{D}_{n} = (-1)^{n+1} \mathrm{P}_{n} \circ \Delta_{n} \circ \mathrm{I}_{n}\) ；于是我们有

\[
\mathrm{Z} \left(\operatorname{Homgr} _ {\mathrm{A}} \left(\mathrm{C}, \mathrm{C} ^ {\prime}\right)\right) = \operatorname{Ker} (\mathrm{P} \circ \Delta \circ \mathrm{I}) = \operatorname{Ker} \mathrm{I} = \operatorname{Im} \mathrm{P},
\]

\[
\mathrm{B} \left(\operatorname{Homgr} _ {\mathrm{A}} \left(\mathrm{C}, \mathrm{C} ^ {\prime}\right)\right) = \operatorname{Im} (\mathrm{P} \circ \Delta \circ \mathrm{I}) = \mathrm{P} (\operatorname{Im} \Delta) = \mathrm{P} (\text { Ker   J });
\]

由此得到一个同构 \(\varphi : H(\text{Homgr}_A(C, C')) \to \text{Homgr}_A(H(C), C')\) ，使得如果 \(a \in \text{Homgr}_A(C/B(C), C')\) ，则在 \(\varphi\) 下 \(P(a)\) 的类的像就是 \(J(a)\) ；立即验证可得 \(\varphi\) 与典范同态 \(\lambda\) 等同.

推论2. --- 假设 B(C) 和 H(C) 是投射的（分别地，B \(_{n}\) (C') 和 H \(_{n}\) (C') 对每个 n 都是内射的)。则 λ(C, C') 是双射。

事实上，C（相应地 C′）此时按 X 第 35 页例 4 是分裂的，我们应用推论1。

推论3.——设 M 为投射（相应地，内射）A-模。对于 A-模的每个复形 C 及每个整数 n，典范同态

\[
\mathrm{H} ^ {n} (\operatorname{Homgr} _ {\mathbf {A}} (\mathrm{M}, \mathrm{C})) \rightarrow \operatorname{Hom} _ {\mathbf {A}} (\mathrm{M}, \mathrm{H} ^ {n} (\mathrm{C}))
\]

\[
\left(\text { resp. } \mathrm{H} ^ {n} (\operatorname{Homgr} _ {\mathbf {A}} (\mathbf {C}, \mathbf {M})) \rightarrow \operatorname{Hom} _ {\mathbf {A}} \left(\mathrm{H} _ {n} (\mathbf {C}), \mathbf {M}\right)\right) \text { is   bijective. }
\]

引理1. --- a) 若 C 或 C′ 右有界，且 C 是投射的、H(C′) = 0，则 H(Homgr\_A(C, C')) = 0。

\begin{enumerate}
\def\labelenumi{\alph{enumi})}
\setcounter{enumi}{1}
\tightlist
\item
  若 C 或 C' 左有界，且 C' 是内射的、H(C)=0，则 H(Homgr\_A(C, C')) = 0。
\end{enumerate}

设 \(f \in Z_n(\text{Homgr}_A(C, C'))\) ；\(f\) 因而是从 \(C\) 到 \(C'(n)\) 的复形态射；在情形 \(a\) （相应地， \(b\) ）， \(f_m\) 当 \(m\) 足够小（分别地，足够大）时为零。根据 \(X\) ，第 47 页，命题 1， \(f\) 于是同伦于零，因此属于 \(B_n(\text{Homgr}_A(C, C'))\) ，由此得出结论。

命题4. --- 设 \(u: \mathbb{C}' \to \mathbb{C}\) 是复形的拟同构， \(\mathbb{P}\) 一个投射复形， \(\mathbb{E}\) 一个内射复形。

\begin{enumerate}
\def\labelenumi{\alph{enumi})}
\tightlist
\item
  若 P 上有界，或 C 与 C' 上有界，则
\end{enumerate}

\[
\operatorname{Homgr} _ {\mathbf {A}} (1, u): \operatorname{Homgr} _ {\mathbf {A}} (\mathrm{P}, \mathrm{C} ^ {\prime}) \rightarrow \operatorname{Homgr} _ {\mathbf {A}} (\mathrm{P}, \mathrm{C})
\]

是拟同构。

\begin{enumerate}
\def\labelenumi{\alph{enumi})}
\setcounter{enumi}{1}
\tightlist
\item
  若 E 下有界，或 C 与 C' 下有界，则
\end{enumerate}

\[
\operatorname{Homgr} _ {\mathbf {A}} (u, 1): \operatorname{Homgr} _ {\mathbf {A}} (\mathrm{C}, \mathrm{E}) \rightarrow \operatorname{Homgr} _ {\mathbf {A}} (\mathrm{C} ^ {\prime}, \mathrm{E})
\]

是拟同构。

首先假设 u 是单射的，并设 \(C'' = \text{Coker } u\) 。由于 u 是一个拟同构， \(C''\) 的同调为零。另一方面，我们有正合序列（命题 2）

\[
0 \rightarrow \operatorname{Homgr} _ {\mathrm{A}} (\mathrm{P}, \mathrm{C} ^ {\prime}) \xrightarrow {\operatorname{Homgr} (1 , u)} \operatorname{Homgr} _ {\mathrm{A}} (\mathrm{P}, \mathrm{C}) \rightarrow \operatorname{Homgr} _ {\mathrm{A}} (\mathrm{P}, \mathrm{C} ^ {\prime \prime}) \rightarrow 0
\]

\[
0 \rightarrow \operatorname{Homgr} _ {\mathbf {A}} \left(\mathrm{C} ^ {\prime \prime}, \mathrm{E}\right)\rightarrow \operatorname{Homgr} _ {\mathbf {A}} (\mathrm{C}, \mathrm{E}) \xrightarrow {\operatorname{Homgr} (u , 1)} \operatorname{Homgr} _ {\mathbf {A}} \left(\mathrm{C} ^ {\prime}, \mathrm{E}\right)\rightarrow 0
\]

由引理1，Homgr \(_{A}\) (P, C'') 在情形 \(a\) ) 下同调为零，Homgr \(_{A}\) (C'', E) 在情形 \(b\) ) 下同调为零，由此得出结论。

在一般情况下，存在（X，第 38 页，命题 7 的推论）一个复形 \(\tilde{\mathbf{C}}'\) ，它当 C 和 C' 有界时右（相应地，左）有界，一个单射态射 \(\tilde{u}: \mathbf{C}' \to \tilde{\mathbf{C}}'\) 以及一个同伦等价 \(\beta: \tilde{\mathbf{C}}' \to \mathbf{C}\) ，使得 \(u = \beta \circ \tilde{u}\) 。于是 \(\tilde{u}\) 是一个拟同构（X，第 34 页，命题 5 的推论）；根据前述，Homgr \(_{\mathbf{A}}\) ( \(1_{\mathbf{P}}, \tilde{u}\) ）（相应地，Homgr \(_{\mathbf{A}}\) ( \(\tilde{u}, 1_{\mathbf{E}}\) ））在情形 \(a\) )（相应地， \(b\) ）中是一个拟同构。另一方面，Homgr \(_{\mathbf{A}}\) ( \(1_{\mathbf{P}}, \beta\) ）（相应地，Homgr \(_{\mathbf{A}}\) ( \(\beta, 1_{\mathbf{E}}\) )）是一个同伦等价（命题3）：于是 Homgr \(_{\mathbf{A}}\) ( \(1_{\mathbf{P}}, u\) ）（相应地，Homgr \(_{\mathbf{A}}\) ( \(u, 1_{\mathbf{E}}\) )）由两个拟同构复合而成，因此是一个拟同构。

